%%% FIELD ambient-expansion-statement
For the fixed bounded class \(\Theta\), uniformly over all triangular arrays \(\theta_n\in\Theta\), write the observation in environment \(e\) as \(W_{e,i}=(T_{e,i},Y_{e,i})\), put \(C^+=(C^\top C)^{-1}C^\top\), and put \(h_k=m_k-\beta c_k\).  There are centered first-projection arrays \(\psi^C_{e,n}(W_{e,i})\), \(\psi^c_{e,k,n}(W_{e,i})\), and \(\psi^m_{e,k,n}(W_{e,i})\) such that, simultaneously for \(d+2\leq k\leq R_{\max}\),
\[
 \sqrt n\,\widehat Q(\widehat m_k-\beta\widehat c_k)
 =Z_{n,k}+o_{\mathbb P}(1),
 \qquad
 Z_{n,k}=\frac1{\sqrt n}\sum_{e\in\mathcal E}\sum_{i=1}^{n_e}\pi_e^{-1}\zeta_{e,k,n}(W_{e,i}),
\]
uniformly over \(\Theta\), where
\[
 \zeta_{e,k,n}
 =Q_A\bigl(\psi^m_{e,k,n}-\beta\psi^c_{e,k,n}
       -\psi^C_{e,n}C^+h_k\bigr).
\]
The stacked vector \(Z_n=(Z_{n,k}:d+2\leq k\leq R_{\max})\) satisfies
\[
 \sup_{\theta\in\Theta}
 d_{\mathrm{BL}}\!\left(\mathcal L_\theta(Z_n),N(0,\Sigma_{n,\theta})\right)\longrightarrow0,
 \qquad
 \Sigma_{n,\theta}=\sum_{e\in\mathcal E}\pi_e^{-1}
 \mathbb E_\theta\bigl[\zeta_{e,n}\zeta_{e,n}^{\top}\bigr].
\]
This conclusion does not require a nonzero projected cumulant flag and allows \(\Sigma_{n,\theta}\) to be singular, to change rank along a triangular array, or to equal zero.  The term \(-Q_A\psi^C_{e,n}C^+h_k\) is the same-sample estimated-\(A\) correction.
%%% FIELD ambient-expansion-proof
Fix an arbitrary triangular array \(\theta_n\in\Theta\), and suppress its index on population quantities.  Every bound below depends only on the fixed envelopes defining \(\Theta\), so taking a supremum over the array is legitimate.  Under the stipulated stratified sampling law, observations are independent across strata and, in stratum \(e\), \(W_{e,1},\ldots,W_{e,n_e}\) are identically distributed with law \(P_e\).  Balanced shares give
\[
 n_e=n\pi_e\geq n\underline\pi. \tag{1}
\]
Thus, for all sufficiently large \(n\), every stratum has at least \(R_{\max}\) observations and the small-sample fallback in the cumulant-SVD construction is inactive.

First we derive the observation-moment envelope.  The coefficient restrictions give
\[
 |\alpha|\leq a_1,
 \qquad |\beta|\leq B,
 \qquad |\gamma|=|\alpha|\,|\gamma/\alpha|\leq a_1q_1. \tag{2}
\]
The structural equations are
\[
 T_e=\alpha\epsilon_{u,e}+\epsilon_{t,e},
 \qquad
 Y_e=(\beta\alpha+\gamma)\epsilon_{u,e}
      +\beta\epsilon_{t,e}+\epsilon_y. \tag{3}
\]
Put \(s_0=2R_{\max}+\delta\).  Since \(s_0>1\), convexity yields
\[
 |x_1+\cdots+x_j|^{s_0}
 \leq j^{s_0-1}\sum_{a=1}^j|x_a|^{s_0}. \tag{4}
\]
Applying (4) to (3), then using (2) and the source moment envelope, produces a finite constant \(K_W\) such that
\[
 \sup_{n}\max_{e\in\mathcal E}
 \mathbb E_{\theta_n}\|W_{e,1}\|^{2R_{\max}+\delta}
 \leq K_W. \tag{5}
\]
The moment--cumulant polynomial consequently bounds, uniformly, every population moment and cumulant of \((T_e,Y_e)\) through order \(R_{\max}\).  In particular, \(C\), \(c_k\), \(m_k\), and \(h_k=m_k-\beta c_k\) have uniformly bounded norms for \(2\leq k\leq R_{\max}\).

We next derive the linear representations of the sample cumulants.  For random variables \(X_1,\ldots,X_k\), the defining moment--cumulant polynomial is
\[
 \kappa_k(X_1,\ldots,X_k)
 =\sum_{\mathcal P\in\Pi_k}
 (|\mathcal P|-1)!(-1)^{|\mathcal P|-1}
 \prod_{D\in\mathcal P}
 \mathbb E\!\left[\prod_{j\in D}X_j\right], \tag{6}
\]
where \(\Pi_k\) is the finite collection of partitions of \(\{1,\ldots,k\}\).  For a partition with \(s\) blocks, average the corresponding block monomials over \(s\) distinct observations and symmetrize.  This is an unbiased symmetric \(U\)-statistic of order \(s\).  Hence each sample cumulant stipulated by the estimator is a fixed finite linear combination of symmetric \(U\)-statistics of orders at most \(k\), with kernels of total polynomial degree at most \(k\).

Consider one order-\(s\) constituent \(U_{n_e,s}(g)\).  Successively subtracting conditional expectations gives canonical centered kernels \(g_j\), \(1\leq j\leq s\), and the exact identity
\[
 U_{n_e,s}(g)-\mathbb Eg
 =\frac1{n_e}\sum_{i=1}^{n_e}\psi_g(W_{e,i})
  +\sum_{j=2}^s {s\choose j}U_{n_e,j}(g_j), \tag{7}
\]
where
\[
 \psi_g(w)=s\{\mathbb E[g(w,W_2,\ldots,W_s)]-\mathbb Eg\},
 \qquad \mathbb E\psi_g(W_{e,1})=0. \tag{8}
\]
If two distinct unordered \(j\)-tuples index summands in \(U_{n_e,j}(g_j)\), one tuple contains an index absent from the other.  Conditioning on all remaining variables and using degeneracy in that argument makes their covariance zero.  Therefore
\[
 \mathbb E\!\left[U_{n_e,j}(g_j)^2\right]
 ={n_e\choose j}^{-1}\mathbb E[g_j^2]
 \leq K n_e^{-j}. \tag{9}
\]
The last inequality is uniform: the squared kernel has degree at most \(2R_{\max}\), so (5), H\"older's inequality, and the finite number of cumulant partitions bound its expectation.  Equations (7)--(9) show that the sum of all projections of order at least two is \(O_{L^2}(n_e^{-1})\), uniformly.

Apply the environment-versus-baseline contrast signs and stack the first projections (8).  This defines centered arrays \(\psi^C_{e,n}\), \(\psi^c_{e,k,n}\), and \(\psi^m_{e,k,n}\).  Because \(n_e=n\pi_e\), (1), (7), and (9) imply, simultaneously over the finitely many relevant \(k\),
\[
 \begin{aligned}
 \sqrt n(\widehat C-C)
 &=\frac1{\sqrt n}\sum_{e\in\mathcal E}\sum_{i=1}^{n_e}
     \pi_e^{-1}\psi^C_{e,n}(W_{e,i})+o_{L^2}(1),\\
 \sqrt n(\widehat c_k-c_k)
 &=\frac1{\sqrt n}\sum_{e\in\mathcal E}\sum_{i=1}^{n_e}
     \pi_e^{-1}\psi^c_{e,k,n}(W_{e,i})+o_{L^2}(1),\\
 \sqrt n(\widehat m_k-m_k)
 &=\frac1{\sqrt n}\sum_{e\in\mathcal E}\sum_{i=1}^{n_e}
     \pi_e^{-1}\psi^m_{e,k,n}(W_{e,i})+o_{L^2}(1).
 \end{aligned} \tag{10}
\]
Indeed, multiplying an \(O_{L^2}(n_e^{-1})\) remainder by \(\sqrt n\) gives \(O(n^{-1/2})\) by (1).

Let
\[
 \nu=\min\{1,\delta/R_{\max}\}>0. \tag{11}
\]
Each first projection in (10) is a finite sum of polynomials of degree at most \(R_{\max}\) in one observation, with coefficients formed from uniformly bounded lower-order moments.  Conditional Jensen and (5) therefore give
\[
 \sup_{n,e,\theta_n\in\Theta}\mathbb E_{\theta_n}\!\left[
 \|\psi^C_{e,n}\|^{2+\nu}
 +\sum_{k=2}^{R_{\max}}
 \{\|\psi^c_{e,k,n}\|^{2+\nu}
   +\|\psi^m_{e,k,n}\|^{2+\nu}\}\right]<\infty, \tag{12}
\]
because \((2+\nu)R_{\max}\leq2R_{\max}+\delta\).

We now linearize the generated projector.  Let \(P_C=CC^+\), so \(Q_A=I_{E-1}-P_C\).  The calibration gap licenses every division below and yields
\[
 \|C^+\|_{\mathrm{op}}=\sigma_d(C)^{-1}\leq\kappa_0^{-1}. \tag{13}
\]
For \(\|H\|_{\mathrm{op}}\leq\kappa_0/2\),
\[
 \sigma_d(C+H)\geq\sigma_d(C)-\|H\|_{\mathrm{op}}
 \geq\kappa_0/2. \tag{14}
\]
Writing \(G=C^\top C\) and \(D=C^\top H+H^\top C+H^\top H\), use the exact inverse identity
\[
 (G+D)^{-1}=G^{-1}-G^{-1}DG^{-1}
 +G^{-1}DG^{-1}D(G+D)^{-1}. \tag{15}
\]
Substitution into \(P_{C+H}=(C+H)\{(C+H)^\top(C+H)\}^{-1}(C+H)^\top\), followed by collection of the terms linear in \(H\), gives
\[
 P_{C+H}-P_C
 =Q_AHC^+ +(Q_AHC^+)^\top+R_C(H),
 \qquad \|R_C(H)\|_{\mathrm{op}}\leq K_P\|H\|_{\mathrm{op}}^2. \tag{16}
\]
The constant is uniform because (5)--(6) bound \(\|C\|\), (13) bounds \(G^{-1}\), and (14) bounds the final inverse in (15).

Taking \(H=\widehat C-C\), equations (10) and (12) give \(\|H\|=O_{\mathbb P}(n^{-1/2})\) uniformly.  Thus (16) applies with probability tending to one uniformly, and its remainder is \(O_{\mathbb P}(n^{-1})\).  On this event the rank-\(d\) SVD projector is the column-space projector of \(\widehat C\).  Hence
\[
 \widehat Q-Q_A
 =-Q_A(\widehat C-C)C^+
  -\{Q_A(\widehat C-C)C^+\}^\top
  +o_{\mathbb P}(n^{-1/2}) \tag{17}
\]
uniformly over \(\Theta\).

The cumulant-cancellation lemma applies to every \(\theta_n\in\Theta\) and gives
\[
 h_k=m_k-\beta c_k=\alpha^{k-1}\gamma u_k\in A=\operatorname{col}(C),
 \qquad Q_Ah_k=0. \tag{18}
\]
In particular,
\[
 \{Q_A(\widehat C-C)C^+\}^\top h_k
 =C^{+\top}(\widehat C-C)^\top Q_Ah_k=0. \tag{19}
\]
Let \(\widehat h_k=\widehat m_k-\beta\widehat c_k\).  Expanding the product, applying (10) and (17)--(19), and using the uniform bounds on \(\beta\), \(h_k\), and \(C^+\), gives
\[
 \begin{aligned}
 \sqrt n\,\widehat Q\widehat h_k
 &=\sqrt n\,Q_A(\widehat h_k-h_k)
   +\sqrt n(\widehat Q-Q_A)h_k
   +\sqrt n(\widehat Q-Q_A)(\widehat h_k-h_k)\\
 &=\sqrt n\,Q_A\{(\widehat m_k-m_k)
       -\beta(\widehat c_k-c_k)\}
   -\sqrt n\,Q_A(\widehat C-C)C^+h_k+o_{\mathbb P}(1). 
 \end{aligned} \tag{20}
\]
The product in the first line is \(O_{\mathbb P}(n^{-1/2})\), uniformly, because each factor before multiplication by \(\sqrt n\) is \(O_{\mathbb P}(n^{-1/2})\).  Substituting (10) into (20) proves
\[
 \sqrt n\,\widehat Q(\widehat m_k-\beta\widehat c_k)
 =\frac1{\sqrt n}\sum_{e\in\mathcal E}\sum_{i=1}^{n_e}
   \pi_e^{-1}\zeta_{e,k,n}(W_{e,i})+o_{\mathbb P}(1), \tag{21}
\]
with
\[
 \zeta_{e,k,n}
 =Q_A\{\psi^m_{e,k,n}-\beta\psi^c_{e,k,n}
          -\psi^C_{e,n}C^+h_k\}. \tag{22}
\]
This proves both the sign and the form of the same-sample estimated-plane correction.  No step used \(r^\star\), \(\lambda_n\), or nonvanishing of a projected cumulant.

It remains to establish the Gaussian approximation without a covariance-rank condition.  Stack (22) over \(d+2\leq k\leq R_{\max}\).  Equations (12)--(13) and the uniform bound on \(h_k\) imply
\[
 \sup_{n,e,\theta_n\in\Theta}
 \mathbb E_{\theta_n}\|\zeta_{e,n}\|^{2+\nu}<\infty. \tag{23}
\]
Index independent summands by \(a=(e,i)\) and put
\[
 X_a=n^{-1/2}\pi_e^{-1}\zeta_{e,n}(W_{e,i}),
 \qquad L_n=\sum_a\mathbb E\|X_a\|^{2+\nu}.
\]
By (1) and (23),
\[
 L_n\leq K n^{-1-\nu/2}\sum_{e\in\mathcal E}n_e\pi_e^{-(2+\nu)}
 \leq K'n^{-\nu/2}. \tag{24}
\]
For each \(a\), let \(G_a\) be an independent centered Gaussian vector with the covariance of \(X_a\).  Fixed dimensionality and Lyapunov's inequality imply
\[
 \mathbb E\|G_a\|^{2+\nu}
 \leq K\{\mathbb E\|X_a\|^2\}^{1+\nu/2}
 \leq K\mathbb E\|X_a\|^{2+\nu}. \tag{25}
\]
For a function \(f\) bounded by one and one-Lipschitz, let \(f_\varepsilon(x)=\mathbb E[f(x+\varepsilon Z_0)]\), where \(Z_0\) is standard Gaussian in the fixed stacked dimension.  Translation coupling gives \(\|f_\varepsilon-f\|_\infty\leq K\varepsilon\).  Differentiation of the Gaussian convolution gives
\[
 \sup_x\|D^2f_\varepsilon(x)\|\leq K\varepsilon^{-1},
 \qquad
 \sup_x\|D^3f_\varepsilon(x)\|\leq K\varepsilon^{-2}.
\]
Interpolating these two bounds yields
\[
 \|D^2f_\varepsilon(x)-D^2f_\varepsilon(y)\|
 \leq K\varepsilon^{-(1+\nu)}\|x-y\|^\nu. \tag{26}
\]
The integral second-order Taylor remainder is consequently bounded by
\[
 K\varepsilon^{-(1+\nu)}\|v\|^{2+\nu}. \tag{27}
\]
Replace the \(X_a\)'s successively by the \(G_a\)'s.  At each replacement, the constant, linear, and quadratic Taylor terms have identical expectations because the paired variables have the same mean and covariance.  Equations (25) and (27) therefore yield
\[
 \left|\mathbb Ef_\varepsilon\!\left(\sum_aX_a\right)
       -\mathbb Ef_\varepsilon\!\left(\sum_aG_a\right)\right|
 \leq K\varepsilon^{-(1+\nu)}L_n. \tag{28}
\]
Combining smoothing and replacement,
\[
 d_{\mathrm{BL}}\!\left(\mathcal L\!\left(\sum_aX_a\right),
                         \mathcal L\!\left(\sum_aG_a\right)\right)
 \leq K\{\varepsilon+\varepsilon^{-(1+\nu)}L_n\}. \tag{29}
\]
When \(L_n>0\), take \(\varepsilon=L_n^{1/(2+\nu)}\); when \(L_n=0\), both sums vanish almost surely.  Equations (24) and (29) prove uniform convergence to zero.  Finally,
\[
 \operatorname{Cov}\!\left(\sum_aG_a\right)
 =\frac1n\sum_{e\in\mathcal E}n_e\pi_e^{-2}
   \mathbb E[\zeta_{e,n}\zeta_{e,n}^\top]
 =\sum_{e\in\mathcal E}\pi_e^{-1}
   \mathbb E[\zeta_{e,n}\zeta_{e,n}^\top]
 =\Sigma_{n,\theta}. \tag{30}
\]
Thus the comparison law is \(N(0,\Sigma_{n,\theta})\).  The construction of the independent Gaussian summands is valid for every positive semidefinite covariance.  If \(\Sigma_{n,\theta}\) is singular, the law is supported on its image; if it is zero, the law is a point mass.  This proves the ambient-class expansion and all its uniformity assertions.
%%% FIELD ambient-randomized-statement
Put
\[
 \nu=\min\{1,\delta/R_{\max}\},
 \qquad
 \varrho_{\mathrm{AR}}=\frac{\nu}{8+4\nu}.
\]
Let \(0<\chi<\varrho_{\mathrm{AR}}\) and \(m_n\to\infty\) be deterministic.  The randomized-jitter Gaussian-multiplier set \(\mathcal C_n^{\mathrm{RJ}}\) is implementable from the stratified observations and the specified auxiliary standard normal draws, and
\[
 \lim_{n\to\infty}
 \sup_{\theta\in\Theta}
 \left|
 \Pr_\theta\{\beta\in\mathcal C_n^{\mathrm{RJ}}\}
 -(1-\alpha_0)
 \right|=0.
\]
The probability is unconditional over the observations, all Gaussian multipliers, the observed-statistic jitter, and all Monte Carlo jitters.  This ambient-class conclusion includes parameters with no projected cumulant breakout.  It allows the same-sample estimated-\(A\) influence covariance \(\Sigma_{n,\theta}\) to be singular, to change rank along a triangular array, or to equal zero, and it uses neither a covariance inverse nor a positive-eigenvalue bound.  Since \(\tau_n=n^{-\chi}\to0\), the external randomization vanishes on the root-\(n\) projected-moment scale underlying \(\mathcal C_n\).  No expected-length or diameter optimality is asserted for \(\mathcal C_n^{\mathrm{RJ}}\).

There is also a precise atomic limitation for deterministic single-threshold calibration.  If a reference statistic has distribution function \(F\), has an atom at a proposed cutoff \(x_0\), and
\[
 F(x_0-)<1-\alpha_0<F(x_0),
\]
then neither the deterministic rule accepting below \(x_0\) nor the deterministic rule accepting at or below \(x_0\) has acceptance probability \(1-\alpha_0\).  This limitation concerns only those deterministic single-threshold conventions at an atom; it is not an impossibility claim about every procedure without external randomization.
%%% FIELD ambient-randomized-proof
Write \(\mathcal K=\{d+2,\ldots,R_{\max}\}\) and use the block maximum norm in the randomized-jitter construction.  The dimension \(D_{\mathrm{AR}}=(E-1)(R_{\max}-d-1)\) is fixed.  Throughout the proof, constants depend only on the fixed coefficient, moment, share, and calibration-gap envelopes and not on \(\theta\in\Theta\).

Balanced shares imply \(n_e\geq\underline\pi n\), so every \(n_e\) is eventually at least \(R_{\max}\).  Thus the construction is defined for every sample size by its fallback and eventually uses the stated cumulant \(U\)-statistics in every stratum.  At the true candidate \(b=\beta\), the ambient generated-projector expansion gives
\[
 H_n(\beta)=Z_n+R_n,
 \qquad
 Z_n=\frac1{\sqrt n}\sum_{e\in\mathcal E}\sum_{i=1}^{n_e}
 \pi_e^{-1}\zeta_{e,n}(W_{e,i}), \tag{1}
\]
where \(R_n=o_{\mathbb P}(1)\) uniformly over \(\Theta\), and
\[
 \Sigma_{n,\theta}=\sum_{e\in\mathcal E}\pi_e^{-1}
 \mathbb E_\theta[\zeta_{e,n}\zeta_{e,n}^{\top}]. \tag{2}
\]
That expansion uses the exact identity \(Q_A(m_k-\beta c_k)=0\), so neither (1) nor any argument below requires a breakout.  With
\[
 \nu=\min\{1,\delta/R_{\max}\}>0,
\]
its construction and the moment calculation in its proof give
\[
 \sup_{n,e,\theta\in\Theta}
 \mathbb E_\theta\|\zeta_{e,n}\|^{2+\nu}\leq K. \tag{3}
\]
The required source moment order is at most \((2+\nu)R_{\max}\leq2R_{\max}+\delta\), and every occurrence of \(C^+\) is bounded because \(\sigma_d(C)\geq\kappa_0>0\).

We first record quantitative Gaussian and expansion-remainder bounds.  Put
\[
 X_{e,i,n}=n^{-1/2}\pi_e^{-1}\zeta_{e,n}(W_{e,i}).
\]
Equation (3) and balanced shares imply
\[
 \sum_{e\in\mathcal E}\sum_{i=1}^{n_e}
 \mathbb E\|X_{e,i,n}\|^{2+\nu}
 \leq K n^{-\nu/2}. \tag{4}
\]
The smoothing and Lindeberg replacement calculation in the ambient expansion proof is quantitative: smoothing at bandwidth \(s\) bounds the bounded-Lipschitz distance by
\[
 K\{s+n^{-\nu/2}s^{-1-\nu}\}.
\]
Taking \(s=n^{-\nu/(4+2\nu)}\) yields
\[
 \sup_{\theta\in\Theta}
 d_{\mathrm{BL}}\bigl(\mathcal L_\theta(Z_n),N(0,\Sigma_{n,\theta})\bigr)
 \leq K n^{-\nu/(4+2\nu)}. \tag{5}
\]
This replacement compares covariance matrices directly and does not divide by an eigenvalue.

For later smoothing we also need a rate for \(R_n\), not merely convergence in probability.  For a canonical Hoeffding projection of order \(j\geq2\), the distinct-index covariance cancellation proved in the expansion gives mean square \(O(n_e^{-j})\) before root-\(n\) scaling, hence at most \(O(n^{-1})\) afterwards.  The spectral-projector remainder is bounded, on \(\{\|\widehat C-C\|\leq\kappa_0/2\}\), by \(K\sqrt n\|\widehat C-C\|^2\).  The first-projection variance calculation gives
\[
 \mathbb E\|\widehat C-C\|^2\leq K/n,
 \qquad
 \Pr\{\|\widehat C-C\|>\kappa_0/2\}\leq K/n. \tag{6}
\]
Thus
\[
 \mathbb E\!\left[1\wedge\{\sqrt n\|\widehat C-C\|^2\}\right]
 \leq K n^{-1/2}.
\]
Products of two first-order estimation errors have the same truncated-mean order by Cauchy--Schwarz and their \(O(n^{-1})\) mean squared errors.  Combining these facts gives
\[
 \sup_{\theta\in\Theta}\mathbb E_\theta[1\wedge\|R_n\|]
 \leq K n^{-1/2}. \tag{7}
\]
A bounded one-Lipschitz test function changes by at most \(1\wedge\|R_n\|\).  Since \(\nu\leq1\), (5)--(7) imply
\[
 \sup_{\theta\in\Theta}
 d_{\mathrm{BL}}\bigl(\mathcal L_\theta(H_n(\beta)),
 N(0,\Sigma_{n,\theta})\bigr)
 \leq K n^{-\nu/(4+2\nu)}. \tag{8}
\]

We next prove the quantitative estimated-influence covariance bound.  This step is necessary because convergence in probability alone is insufficient after division by the vanishing jitter scale.  For an order-\(s\) constituent kernel \(g\), define its population first projection by
\[
 \psi_{g,e}(w)=s\{\mathbb E[g(w,W_2,\ldots,W_s)]-\mathbb Eg\}.
\]
For \(s\geq2\), conditionally on \(W_{e,i}\), the first average in \(\widehat\psi_{g,e,i}\) is an order-\((s-1)\) \(U\)-statistic based on the remaining \(n_e-1\) observations.  Among two indexing tuples, disjoint tuples are conditionally independent, and the number of overlapping ordered pairs is \(O(n_e^{2s-3})\), compared with \(O(n_e^{2s-2})\) total pairs.  The kernel second moment is uniformly bounded by the source moment envelope and the structural coefficient bounds.  The same argument applies to \(U_{e,s}(g)\).  For \(s=1\), the conditional-estimation term is absent.  Therefore, after stacking the finitely many kernels and contrast signs,
\[
 \mathbb E R_n^\psi\leq K/n,
 \qquad
 R_n^\psi=\frac1n\sum_{e\in\mathcal E}\sum_{i=1}^{n_e}\pi_e^{-2}
 \|\widehat\psi_{e,i}-\psi_{e,n}(W_{e,i})\|^2. \tag{9}
\]
Conditional Jensen applied to the leave-one-out averages and ordinary Jensen applied to the \(U\)-statistics give uniform \((2+\nu)\)-moment bounds for \(\widehat\psi_{e,i}\) and \(\psi_{e,n}(W_{e,i})\), because their polynomial degree is at most \(R_{\max}\).

Put \(h_k=m_k-\beta c_k\), \(\widehat h_k=\widehat m_k-\beta\widehat c_k\), and
\[
 L_n^2=\|\widehat Q-Q_A\|_{\mathrm{op}}^2
 +\|\widehat C^\dagger-C^+\|_{\mathrm{op}}^2
 +\sum_{k\in\mathcal K}\|\widehat h_k-h_k\|^2. \tag{10}
\]
On \(\{\|\widehat C-C\|\leq\kappa_0/2\}\), the inverse identity used for the projector gives \(\|\widehat C^\dagger-C^+\|\leq K\|\widehat C-C\|\).  Off this event, the threshold in the definition gives \(\|\widehat C^\dagger\|\leq2/\kappa_0\), while \(\|C^+\|\leq1/\kappa_0\), and (6) controls its probability.  The projector perturbation bound and the sample-cumulant mean squared errors consequently yield
\[
 \mathbb E L_n^2\leq K/n,
 \qquad
 \mathbb E L_n\leq K n^{-1/2}. \tag{11}
\]
Let \(a=\nu/2\) and define
\[
 Y_n=1+\frac1n\sum_{e\in\mathcal E}\sum_{i=1}^{n_e}\pi_e^{-2}
 \{\|\psi_{e,n}(W_{e,i})\|^2+\|\widehat\psi_{e,i}\|^2\}.
\]
Convexity and the \((2+\nu)\)-moment bounds imply
\[
 \mathbb E Y_n^{1+a}\leq K. \tag{12}
\]
Expanding \(\widehat\zeta_{e,i}(\beta)-\zeta_{e,n}(W_{e,i})\), adding and subtracting one estimated factor at a time, and using the bounds on \(C^+\), \(\widehat C^\dagger\), and \(h_k\), gives the deterministic inequality
\[
 B_n:=\frac1n\sum_{e\in\mathcal E}\sum_{i=1}^{n_e}\pi_e^{-2}
 \|\widehat\zeta_{e,i}(\beta)-\zeta_{e,n}(W_{e,i})\|^2
 \leq K\{R_n^\psi+L_n^2Y_n\}. \tag{13}
\]
No independence between \(L_n\) and \(Y_n\) is used.

Define
\[
 \widetilde\Sigma_n=\frac1n\sum_{e\in\mathcal E}\sum_{i=1}^{n_e}\pi_e^{-2}
 \zeta_{e,n}(W_{e,i})\zeta_{e,n}(W_{e,i})^\top
\]
and let \(\widehat\Sigma_n(\beta)\) be the analogous empirical covariance built from \(\widehat\zeta_{e,i}(\beta)\).  Empirical Cauchy--Schwarz and (13) give
\[
 \|\widehat\Sigma_n(\beta)-\widetilde\Sigma_n\|_{\mathrm{op}}
 \leq K\{(R_n^\psi Y_n)^{1/2}+L_nY_n+R_n^\psi+L_n^2Y_n\}. \tag{14}
\]
We integrate this inequality without assuming independence of its factors.  For every \(t>0\), (9) and (12) imply
\[
 \mathbb E[1\wedge R_n^\psi Y_n]
 \leq t\mathbb ER_n^\psi+\Pr(Y_n>t)
 \leq K\{t/n+t^{-(1+a)}\}.
\]
Taking \(t=n^{1/(2+a)}\) and then using Jensen gives
\[
 \mathbb E[1\wedge(R_n^\psi Y_n)^{1/2}]
 \leq K n^{-\kappa},
 \qquad
 \kappa=\frac{1+a}{2(2+a)}. \tag{15}
\]
Likewise, (11)--(12) give
\[
 \mathbb E[1\wedge L_nY_n]
 \leq K\{tn^{-1/2}+t^{-(1+a)}\}.
\]
With \(t=n^{1/\{2(2+a)\}}\), this is at most \(K n^{-\kappa}\).  The remaining terms in (14) are faster: \(\mathbb ER_n^\psi\leq K/n\), while choosing \(t=n^{1/(2+a)}\) gives
\[
 \mathbb E[1\wedge L_n^2Y_n]
 \leq K n^{-(1+a)/(2+a)}.
\]
Consequently,
\[
 \sup_{\theta\in\Theta}\mathbb E_\theta\!\left[
 1\wedge\|\widehat\Sigma_n(\beta)-\widetilde\Sigma_n\|_{\mathrm{op}}
 \right]
 \leq K n^{-\kappa}. \tag{16}
\]

For the empirical fluctuation of the true influences, each matrix \(\zeta_{e,n}\zeta_{e,n}^\top\) has a uniformly bounded \((1+a)\)-moment by (3).  Truncate its norm at \(\ell_n=n^{1/\{2(1+a)\}}\).  The truncation bias is at most \(K\ell_n^{-a}\).  Independence across observations, fixed dimension, and balanced shares bound the mean norm of the centered truncated average by its root mean square,
\[
 K n^{-1/2}\ell_n^{(1-a)/2}.
\]
The latter is no larger in order than the bias for \(0<a\leq1/2\).  Moreover,
\[
 \mathbb E\widetilde\Sigma_n
 =\frac1n\sum_e n_e\pi_e^{-2}\mathbb E[\zeta_{e,n}\zeta_{e,n}^\top]
 =\Sigma_{n,\theta}.
\]
It follows that
\[
 \sup_{\theta\in\Theta}\mathbb E_\theta
 \|\widetilde\Sigma_n-\Sigma_{n,\theta}\|_{\mathrm{op}}
 \leq K n^{-a/\{2(1+a)\}}
 =K n^{-\nu/(4+2\nu)}. \tag{17}
\]
By the definition of \(\varrho_{\mathrm{AR}}\),
\[
 2\varrho_{\mathrm{AR}}=\frac{a}{2(1+a)},
 \qquad
 \kappa-2\varrho_{\mathrm{AR}}
 =\frac{1}{2(2+a)(1+a)}>0. \tag{18}
\]
Since \(\tau_n=n^{-\chi}\) and \(\chi<\varrho_{\mathrm{AR}}\), equations (16)--(18) yield the required quantitative replacement statements
\[
 \sup_{\theta\in\Theta}\mathbb E_\theta\!\left[
 1\wedge\|\widehat\Sigma_n(\beta)-\widetilde\Sigma_n\|_{\mathrm{op}}
 \right]=o(\tau_n^2), \tag{19}
\]
\[
 \sup_{\theta\in\Theta}\mathbb E_\theta\!\left[
 1\wedge\|\widehat\Sigma_n(\beta)-\Sigma_{n,\theta}\|_{\mathrm{op}}
 \right]
 \leq K n^{-2\varrho_{\mathrm{AR}}}=o(\tau_n^2). \tag{20}
\]

Conditionally on the observations, \(G_n^{*(j)}(\beta)\) is exactly Gaussian with covariance \(\widehat\Sigma_n(\beta)\).  For positive semidefinite matrices \(S,T\), the integral representation
\[
 S^{1/2}=\frac1\pi\int_0^\infty t^{-1/2}S(S+tI)^{-1}\,dt
\]
and the resolvent identity, with the integral split at \(t=\|S-T\|_{\mathrm{op}}\), give
\[
 \|S^{1/2}-T^{1/2}\|_{\mathrm{op}}
 \leq K\|S-T\|_{\mathrm{op}}^{1/2}. \tag{21}
\]
This derivation remains valid when either matrix is singular.  Couple the conditional multiplier Gaussian and \(N(0,\Sigma_{n,\theta})\) as \(\widehat\Sigma_n(\beta)^{1/2}Z\) and \(\Sigma_{n,\theta}^{1/2}Z\), with \(Z\sim N(0,I_{D_{\mathrm{AR}}})\).  Equations (20)--(21), truncation at one, and Jensen's inequality give
\[
 \sup_{\theta\in\Theta}\mathbb E_\theta
 d_{\mathrm{BL}}\bigl(
 \mathcal L^*(G_n^{*(j)}(\beta)),N(0,\Sigma_{n,\theta})
 \bigr)
 \leq K n^{-\varrho_{\mathrm{AR}}}. \tag{22}
\]

It remains to remove every atomic discontinuity.  Let \(G_{n,\theta}\sim N(0,\Sigma_{n,\theta})\), let \(\xi\sim N(0,1)\) be independent, and define
\[
 F_{n,\theta}^{\circ}(t)
 =\Pr\{\|G_{n,\theta}\|_{\max}+\tau_n\xi\leq t\}.
\]
For each \(t\), the map
\[
 x\longmapsto
 \Phi\!\left(\frac{t-\|x\|_{\max}}{\tau_n}\right)
\]
is bounded by one and has Lipschitz constant at most \((\sqrt{2\pi}\tau_n)^{-1}\).  Hence (8), (22), and \(\chi<\varrho_{\mathrm{AR}}\) imply
\[
 \sup_{\theta\in\Theta,t\in\mathbb R}
 \left|\Pr_\theta\{V_n(\beta)\leq t\}
 -F_{n,\theta}^{\circ}(t)\right|\longrightarrow0, \tag{23}
\]
\[
 \sup_{\theta\in\Theta}\mathbb E_\theta\sup_{t\in\mathbb R}
 \left|\Pr^*\{V_n^{*(j)}(\beta)\leq t\}
 -F_{n,\theta}^{\circ}(t)\right|\longrightarrow0. \tag{24}
\]
Indeed, the respective bounds are \(O(n^{\chi-2\varrho_{\mathrm{AR}}})\) and \(O(n^{\chi-\varrho_{\mathrm{AR}}})\).

The ideal distribution is continuous and strictly increasing even if \(\Sigma_{n,\theta}\) is singular or zero, because it has the everywhere positive density
\[
 t\longmapsto
 \mathbb E\!\left[\tau_n^{-1}\phi\!\left(
 \frac{t-\|G_{n,\theta}\|_{\max}}{\tau_n}\right)\right]. \tag{25}
\]
Let
\[
 p_n^*(\beta)=\Pr^*\{V_n^{*(j)}(\beta)\geq V_n(\beta)\}.
\]
Conditional continuity, supplied by the multiplier jitter, makes the weak inequality immaterial.  Equation (24) shows that \(p_n^*(\beta)\) differs in probability, uniformly over \(\Theta\), from \(1-F_{n,\theta}^{\circ}(V_n(\beta))\).  By (23) and strict continuity of \(F_{n,\theta}^{\circ}\), the latter converges in distribution to the uniform law on \((0,1)\), uniformly: for every \(u\in(0,1)\), applying the inverse of \(F_{n,\theta}^{\circ}\) bounds the distributional error by the left side of (23).  Therefore
\[
 p_n^*(\beta)\ \Longrightarrow\ \operatorname{Unif}(0,1)
 \quad\text{uniformly over }\Theta. \tag{26}
\]

Conditionally on the observations and \(\xi_0\), the \(m_n\) indicators defining \(\widehat p_{n,m_n}(\beta)\) are independent Bernoulli variables with success probability \(p_n^*(\beta)\).  Hence
\[
 \operatorname{Var}\!\left(
 \widehat p_{n,m_n}(\beta)\mid\text{observations},\xi_0
 \right)
 \leq\frac1{4m_n}. \tag{27}
\]
Chebyshev's inequality and \(m_n\to\infty\) make the Monte Carlo error uniformly \(o_{\mathbb P}(1)\).  Since \(\alpha_0\in(0,1/2)\) is a continuity point of the limiting uniform distribution, (26)--(27) give
\[
 \sup_{\theta\in\Theta}
 \left|
 \Pr_\theta\{\widehat p_{n,m_n}(\beta)>\alpha_0\}
 -(1-\alpha_0)
 \right|\longrightarrow0. \tag{28}
\]
The probability in (28) integrates over the observations, the observed jitter, every multiplier, and every Monte Carlo jitter.  By definition, its event is exactly \(\{\beta\in\mathcal C_n^{\mathrm{RJ}}\}\).  This proves ambient-class exact coverage, including at no-breakout laws and under rank-changing or zero influence covariance.  Because \(\tau_n\to0\), both observed and multiplier jitters vanish on the root-\(n\) projected-moment scale.  Nothing in this calibration argument bounds the expected length or diameter of \(\mathcal C_n^{\mathrm{RJ}}\), so no such optimality follows.

Finally, suppose \(F\) has an atom at \(x_0\).  A deterministic strict-threshold rule has acceptance probability \(F(x_0-)\), while a deterministic weak-threshold rule has acceptance probability \(F(x_0)\).  If
\[
 F(x_0-)<1-\alpha_0<F(x_0),
\]
neither probability equals \(1-\alpha_0\).  Randomizing acceptance on the equality event with probability
\[
 \frac{1-\alpha_0-F(x_0-)}{F(x_0)-F(x_0-)}
\]
does attain the target, confirming that the obstruction is exactly to the two deterministic single-threshold conventions and not to every procedure lacking external randomization.
